\documentclass[10pt,letterpaper]{article}

\usepackage[margin=1in]{geometry}
\usepackage{booktabs}
\usepackage{hyperref}

\title{CEPO-Probe Paper Draft Handoff}
\author{Anonymous CVPR Submission}
\date{2026-05-28}

\begin{document}
\maketitle

This prototype draft has been superseded by the submission-style CVPR draft in:

\begin{quote}
\texttt{paper/main.tex}
\end{quote}

The final metric summary for the benchmark-paper pivot is recorded in:

\begin{quote}
\texttt{cepo-probe-benchmark-paper/final\_results.md}
\end{quote}

\section*{Selected Main Groups}

\begin{table}[h]
\centering
\begin{tabular}{ll}
\toprule
Group & Model \\
\midrule
A & Base Instruct \\
B & CEPO Answer-DPO \\
C & CEPO-Dual \\
\bottomrule
\end{tabular}
\end{table}

\section*{Headline Result}

CEPO-Probe shows that short-answer preference tuning and claim-evidence
verification are separable. CEPO Answer-DPO improves ordinary yes/no accuracy
but weakens wrong-evidence rejection, while CEPO-Dual preserves the
short-answer gains and improves explicit evidence consistency from 34.3\% to
70.2\% wrong-evidence rejection. Relation reversals remain the main failure
slice.

\section*{Status}

\begin{itemize}
    \item ClaimEvidence-6K data checks pass with zero train/eval image overlap.
    \item CEPO-Dual is selected as the final evidence-aware baseline.
    \item Main generations were re-scored with the current short-answer and
    evidence-probe scorers on 2026-05-28.
    \item The submission-style manuscript and supplementary material now live
    under \texttt{paper/}.
\end{itemize}

\end{document}
